\documentclass[12pt,a4paper,oneside]{article}
\usepackage{fontspec}
\setmainfont{DejaVu Serif}
\setmonofont{DejaVu Sans Mono}
\usepackage[english,russian]{babel}
\usepackage{geometry}
\geometry{left=2.54cm,right=2.54cm,top=2.54cm,bottom=2.54cm}
\usepackage{amsmath,amssymb,array,longtable,booktabs}
\usepackage{graphicx}
\usepackage{listings}
\usepackage{xcolor}
\usepackage{hyperref}
\usepackage{enumitem}
\usepackage{titlesec}
\titleformat*{\section}{\large\bfseries}
\titleformat*{\subsection}{\normalsize\bfseries}
\setcounter{secnumdepth}{0}
\setlength{\parindent}{0pt}
\setlength{\parskip}{5pt}
\sloppy
\lstset{basicstyle=\footnotesize\ttfamily,frame=single,breaklines=true,columns=fullflexible,numbers=left,numberstyle=\tiny,showstringspaces=false,tabsize=4,keepspaces=true}
\begin{document}
\begin{titlepage}
\begin{center}
\large МОСКОВСКИЙ АВИАЦИОННЫЙ ИНСТИТУТ\\
(НАЦИОНАЛЬНЫЙ ИССЛЕДОВАТЕЛЬСКИЙ УНИВЕРСИТЕТ)
\vspace{20pt}

Институт №8 «Компьютерные науки и прикладная математика»\\
Кафедра 806 «Вычислительная математика и программирование»
\end{center}
\vspace{65pt}
\begin{center}
\large\bfseries Отчет по лабораторной работе №3\\[7pt]
по курсу «Операционные системы»
\end{center}
\vspace{40pt}
\begin{center}
\large\bfseries Тема: «Многопоточная четно-нечетная сортировка Бетчера»
\end{center}
\vfill
\begin{flushright}
\begin{tabular}{rl}
Выполнил(а): & А.\,Е. Пронина \\
Группа: & М8О-203БВ-25 \\
Преподаватель: & Е.\,С. Миронов \\
\end{tabular}
\end{flushright}
\vspace{40pt}
\begin{center}Москва, 2026\end{center}
\end{titlepage}

\section{Условие}
\textbf{Цель работы:} приобретение практических навыков управления потоками операционной системы и обеспечения синхронизации между ними.

\textbf{Задание:} разработать программу для многопоточной обработки данных с использованием стандартных средств Windows или Unix. Максимальное количество потоков, одновременно работающих в программе, задается аргументом командной строки. Необходимо показать потоки средствами операционной системы и исследовать зависимость ускорения и эффективности от размера данных и числа потоков.

\textbf{Вариант 5:} отсортировать массив целых чисел при помощи четно-нечетной сортировки Бетчера.

\section{Метод решения}
\subsection{Общая идея алгоритма}
Используется сортировочная сеть Бетчера (odd-even mergesort network). В отличие от сортировок, которые выбирают следующие сравнения в зависимости от значений элементов, сеть заранее определяет пары \emph{позиций} для сравнения. Компаратор сравнивает два числа и при необходимости меняет их местами, чтобы меньшее значение оказалось слева.

Алгоритм рекурсивно сортирует две половины массива, затем сливает их с помощью четно-нечетной схемы: отдельно формируются сравнения для чередующихся позиций и выполняются завершающие сравнения соседних пар. Полученные операции можно сгруппировать по \emph{уровням}. На одном уровне ни одна позиция не используется двумя компараторами одновременно, а следующий уровень зависит от результата предыдущего.

Для длины $n=2^m$ число уровней классической сети составляет $m(m+1)/2$, общее число сравнений имеет порядок $O(n\log^2 n)$. Глубина сети $O(\log^2 n)$ позволяет выполнять независимые сравнения одного уровня параллельно.

\subsection{Пример построения сети}
Рассмотрим восемь элементов:
\[
[8,3,7,1,6,2,5,4].
\]
Ниже номера позиций указаны \emph{с единицы}. В таблице приведены пары, сравниваемые на каждом уровне, и состояние массива после него.

\begin{center}\small
\begin{tabular}{|c|p{6.3cm}|p{5.0cm}|}
\hline
Уровень & Сравниваемые позиции & Массив после уровня \\ \hline
1 & (1,2), (3,4), (5,6), (7,8) & 3, 8, 1, 7, 2, 6, 4, 5 \\ \hline
2 & (1,3), (2,4), (5,7), (6,8) & 1, 7, 3, 8, 2, 5, 4, 6 \\ \hline
3 & (2,3), (6,7) & 1, 3, 7, 8, 2, 4, 5, 6 \\ \hline
4 & (1,5), (2,6), (3,7), (4,8) & 1, 3, 5, 6, 2, 4, 7, 8 \\ \hline
5 & (3,5), (4,6) & 1, 3, 2, 4, 5, 6, 7, 8 \\ \hline
6 & (2,3), (4,5), (6,7) & 1, 2, 3, 4, 5, 6, 7, 8 \\ \hline
\end{tabular}
\end{center}

На последнем уровне сравниваются позиции (2,3), (4,5) и (6,7), а не произвольные соседние элементы: именно эти пары остаются после рекурсивного четно-нечетного слияния. Например, значения 3 и 2 меняются местами. Внутри одного уровня сравнения независимы и могут быть распределены между потоками.

\subsection{Произвольная длина массива и экономия памяти}
Программа принимает размер массива, не обязательно равный степени двойки. Для построения расписания используется логика полной сети Бетчера; сравнения, которые обращались бы к отсутствующим позициям, пропускаются. Это соответствует мысленному дополнению последовательности фиктивными значениями $+\infty$, но без выделения дополнительного массива для них.

Компараторы не сохраняются в отдельном большом списке. Пара индексов вычисляется непосредственно на текущем уровне сети. Поэтому дополнительная память не увеличивается пропорционально числу всех компараторов. Для входного массива по-прежнему требуется $O(n)$ памяти, а служебные данные потоков занимают $O(P)$, где $P$~--- число участников вычисления.

\subsection{Распределение работы между потоками}
Реализация выполняется на C++ под Ubuntu/WSL с использованием POSIX Threads (\texttt{pthread.h}). Основной поток участвует в вычислениях наряду с рабочими. Если лимит равен $P$, дополнительно создается не более $P-1$ потоков. Для каждого уровня сеть назначает каждому исполнителю собственные позиции исходного массива. Один поток не изменяет элементы, с которыми одновременно работает другой компаратор того же уровня.

Количество рабочих потоков фиксируется на время сортировки: потоки не создаются заново для каждого компаратора. Это уменьшает накладные расходы на их создание и завершение. При $P=1$ сортировка осуществляется только главным потоком.

\subsection{Синхронизация и обработка ошибок}
Для согласования этапов используется \texttt{pthread\_barrier\_wait()}: каждый участник должен завершить сравнения текущего уровня перед переходом к следующему. Без этого один поток мог бы прочитать элемент до того, как его изменит другой поток на предыдущем уровне.

Помимо барьера, используются мьютекс и условная переменная для безопасного начала работы. Созданные потоки ожидают разрешения главного потока. Если создание очередного потока завершилось ошибкой, запуск отменяется: иначе уже созданные потоки могли бы бесконечно ждать недостающего участника барьера.

Главный поток вызывает \texttt{pthread\_join()} для созданных рабочих потоков. После их завершения освобождаются служебная память и объекты синхронизации. Коды возврата функций POSIX должны проверяться; при обнаружении ошибки программа сообщает о ней. В частности, POSIX-функции семейства \texttt{pthread\_*} обычно возвращают код ошибки напрямую, поэтому для диагностики можно применять \texttt{strerror(rc)}.

\section{Описание программы}
Проект организован по модулям:
\begin{itemize}[noitemsep]
\item \texttt{src/main.cpp}~--- обработка аргументов, выделение и заполнение массива, запуск сортировки, вывод;
\item \texttt{src/batcher.h}, \texttt{src/batcher.cpp}~--- формирование и исполнение компараторов сети Бетчера;
\item \texttt{src/threads.h}, \texttt{src/threads.cpp}~--- потоки POSIX, распределение сравнений и синхронизация;
\item \texttt{src/benchmark.h}, \texttt{src/benchmark.cpp}~--- измерение прошедшего времени через \texttt{clock\_gettime()};
\item \texttt{CMakeLists.txt}~--- автоматизированная сборка с подключением библиотеки потоков.
\end{itemize}

Программа запускается командой \texttt{./lab3 N P}, где \texttt{N}~--- число элементов массива, \texttt{P}~--- максимальное число потоков, включая главный. Для небольших массивов печатается исходная и отсортированная последовательность; дополнительно выводятся PID, TID, результат проверки правильности и время сортировки.

\subsection{Ключевые фрагменты реализации}
\textbf{Сравнение двух элементов сети} (сокращенный смысловой фрагмент из \texttt{batcher.cpp}):
\begin{lstlisting}[language=C++]
if (needCompare && arr[i] > arr[i + k]) {
    int temp = arr[i];
    arr[i] = arr[i + k];
    arr[i + k] = temp;
}
\end{lstlisting}
Здесь \texttt{i}~--- индекс первого элемента, \texttt{k}~--- расстояние до второго, а \texttt{needCompare} означает, что эта пара относится к текущему уровню сети. Значение \texttt{needCompare} определяется параметрами уровня и позицией элемента, а не содержимым массива.

\textbf{Барьер после каждого уровня} (\texttt{threads.cpp}):
\begin{lstlisting}[language=C++]
compareStage(arr, n, p, k, id, threadCount);
int rc = pthread_barrier_wait(&barrier);
if (rc != 0 && rc != PTHREAD_BARRIER_SERIAL_THREAD) {
    // Error handling
}
\end{lstlisting}
Смысл фрагмента~--- не допустить начало следующего зависимого уровня, пока все участники не завершили текущий. Возвращаемое значение \texttt{PTHREAD\_BARRIER\_SERIAL\_THREAD} также соответствует успешному ожиданию барьера.

\textbf{Создание и ожидание рабочих потоков} (\texttt{threads.cpp}):
\begin{lstlisting}[language=C++]
int rc = pthread_create(&threads[i], nullptr,
                        worker, &data[i]);
// Check rc before counting the thread as created.

rc = pthread_join(threads[i], nullptr);
// Check rc and release resources after completion.
\end{lstlisting}
Параметры потока передаются через индивидуальную структуру, содержащую его номер и указатель на общие данные. Общие данные включают массив, число участников и объекты синхронизации. Такое разделение не дублирует массив для каждого потока.

\section{Результаты работы}
\subsection{Проверка корректности}
Программа была запущена в Ubuntu через WSL, из терминала VS Code. Для демонстрации выбран массив из 13 целых чисел и лимит 4 потока. Полученный фактический вывод:
\begin{lstlisting}[basicstyle=\scriptsize\ttfamily]
$ ./lab3 13 4
PID: 20186
Array size: 13
Threads: 4
Before: 2946 6373 -5930 6206 1454 9876 -5172
        -8821 -6332 -436 2143 2175 7593
Main TID: 20186
Worker 1 TID: 20187
Worker 2 TID: 20188
Worker 3 TID: 20189
After: -8821 -6332 -5930 -5172 -436 1454
       2143 2175 2946 6206 6373 7593 9876
Sorting: OK
Time: 0.00185605 seconds
\end{lstlisting}
Полученная последовательность отсортирована по неубыванию. Размер 13 не равен степени двойки; следовательно, данный запуск одновременно демонстрирует обработку некратного степени двойки входа. Отображены четыре идентификатора потоков, включая главный. Измеренное время выполнения в этом запуске составило $0{,}00185605$~с.

\subsection{Исследование ускорения и эффективности}
Для исследования производительности проведены запуски на массивах из 1\,000, 10\,000, 100\,000 и 1\,000\,000 элементов при $P=1,2,4,8$ потоках. Для каждого сочетания параметров программа выполнялась три раза. Во всех 48 запусках проверка завершилась сообщением \texttt{Sorting: OK}. В таблице приведены медианы трех измерений времени.

Ускорение и эффективность вычислялись по формулам
\[
S(P)=\frac{T_1}{T_P}, \qquad
E(P)=\frac{S(P)}{P}\cdot100\%,
\]
где $T_1$ и $T_P$~--- медианные времена сортировки массива одинакового размера соответственно с одним и $P$ потоками. При каждом запуске программа генерировала одинаковые исходные данные с фиксированным начальным значением генератора случайных чисел.

\begin{center}
\scriptsize
\setlength{\tabcolsep}{3pt}
\begin{tabular}{|r|r|r|r|r|}
\hline
Размер массива & Потоки $P$ & Время $T_P$, с & Ускорение $S(P)$ & Эффективность $E(P)$ \\ \hline
1\,000 & 1 & 0,000358 & 1,000 & 100,0\% \\ \hline
1\,000 & 2 & 0,003674 & 0,097 & 4,9\% \\ \hline
1\,000 & 4 & 0,007314 & 0,049 & 1,2\% \\ \hline
1\,000 & 8 & 0,014231 & 0,025 & 0,3\% \\ \hline
10\,000 & 1 & 0,006591 & 1,000 & 100,0\% \\ \hline
10\,000 & 2 & 0,011342 & 0,581 & 29,1\% \\ \hline
10\,000 & 4 & 0,015263 & 0,432 & 10,8\% \\ \hline
10\,000 & 8 & 0,025079 & 0,263 & 3,3\% \\ \hline
100\,000 & 1 & 0,100098 & 1,000 & 100,0\% \\ \hline
100\,000 & 2 & 0,070158 & 1,427 & 71,3\% \\ \hline
100\,000 & 4 & 0,055363 & 1,808 & 45,2\% \\ \hline
100\,000 & 8 & 0,055544 & 1,802 & 22,5\% \\ \hline
1\,000\,000 & 1 & 1,322540 & 1,000 & 100,0\% \\ \hline
1\,000\,000 & 2 & 0,775840 & 1,705 & 85,2\% \\ \hline
1\,000\,000 & 4 & 0,625560 & 2,114 & 52,9\% \\ \hline
1\,000\,000 & 8 & 0,374043 & 3,536 & 44,2\% \\ \hline
\end{tabular}
\end{center}

На массивах из 1\,000 и 10\,000 элементов минимальное время показал один поток. Для 100\,000 элементов самый быстрый результат достигнут при четырех потоках ($0{,}055363$~с), а для 1\,000\,000 элементов~--- при восьми ($0{,}374043$~с).

\subsection{Идентификаторы потоков}
В контрольном запуске процесс имел PID 20186. Главный поток имел TID 20186, а рабочие потоки~--- TID 20187, 20188 и 20189. Таким образом, программа создала три дополнительных потока и использовала главный поток как четвертого участника сортировки.

\section{Анализ результатов}
Серия экспериментов показала, что влияние числа потоков существенно зависит от размера массива. При $N=1\,000$ сортировка одним потоком заняла $0{,}000358$~с, тогда как восемью потоками~--- $0{,}014231$~с. При $N=10\,000$ ситуация аналогична: один поток выполнял сортировку за $0{,}006591$~с, а восемь~--- за $0{,}025079$~с. На малых данных затраты на создание потоков и ожидание барьеров превышают выигрыш от параллельного выполнения сравнений.

При $N=100\,000$ многопоточность уже дает ускорение. Время уменьшилось с $0{,}100098$~с на одном потоке до $0{,}055363$~с на четырех, то есть в $1{,}808$ раза. Увеличение числа потоков до восьми не улучшило результат: медианное время составило $0{,}055544$~с. Дополнительные участники увеличивают расходы на синхронизацию и конкурируют за вычислительные ресурсы.

На массиве из $1\,000\,000$ элементов быстрее всего работали восемь потоков: $0{,}374043$~с против $1{,}322540$~с для одного потока. Ускорение составило $3{,}536$, а эффективность~--- $44{,}2\%$. Максимальная эффективность среди многопоточных запусков для этого размера достигнута при двух потоках ($85{,}2\%$), хотя минимальное абсолютное время показали восемь потоков. Таким образом, наилучшее время и наибольшая эффективность использования каждого потока не всегда достигаются при одном и том же значении $P$.

Ускорение ограничивается необходимостью синхронизации после каждого уровня сети, накладными расходами потоков и доступом к общей памяти. При увеличении массива доля собственно вычислений возрастает, поэтому использование нескольких потоков становится выгоднее. Проверки \texttt{Sorting: OK} во всех 48 запусках подтверждают корректность порядка выходных элементов для исследованных размеров и числа потоков.

\section{Выводы}
В лабораторной работе реализована многопоточная четно-нечетная сортировка Бетчера на C++ с использованием POSIX Threads в Ubuntu/WSL. Сортировка принимает массив произвольной длины, вычисляет необходимые пары индексов без хранения всей сети компараторов и выполняет независимые сравнения несколькими потоками.

Было изучено создание потоков через \texttt{pthread\_create()}, ожидание их завершения через \texttt{pthread\_join()} и синхронизация с помощью барьера. Мьютекс и условная переменная обеспечивают согласованный запуск рабочих потоков. Установленный пользователем предел числа потоков включает главный поток, который также участвует в вычислениях.

В ходе испытаний выполнено 48 запусков на четырех размерах массива и четырех значениях числа потоков. Во всех случаях сортировка завершилась корректно. Для небольших массивов эффективнее последовательное выполнение: накладные расходы на многопоточность превосходят пользу от параллельных сравнений. Для массива из 100\,000 элементов лучшее время достигнуто при четырех потоках, а для массива из 1\,000\,000 элементов~--- при восьми. На самом большом массиве получено ускорение $3{,}536$ раза относительно однопоточного выполнения.

Эксперименты показали, что количество потоков следует оценивать вместе с размером данных: увеличение параллелизма ускоряет большие вычисления, но может замедлить обработку небольших массивов из-за стоимости синхронизации и управления потоками.

\section{Список литературы}
\begin{enumerate}[noitemsep]
\item Сеть обменной сортировки со слиянием Бэтчера. Habr. \url{https://habr.com/ru/articles/275889/}.
\end{enumerate}
\end{document}
